\subsection{An affine relation with no prime-field root}
A nonzero constant term can force both the inner polynomial and its
monomial denominator to vanish. The resulting replacement has two
rows, so its complexity must be measured by their total support.

\begin{theorem}[Sparse recognition of an affine relation]
\label{sp:thm:affine-preprocessing}
Fix $p$ and let a homogeneous sparse binary-exponent list generate
$K\subset S=\Fp[x_0,\ldots,x_N]$, with $I=K^{\mathrm{sat}}$ and
$S/I$ geometrically integral. Suppose a nonzero row has a representation
\begin{equation}
 F=U^Q+\sum_{i=0}^s a_iV^{Q-p^i}U^{p^i}+\beta V^Q,
 \qquad Q=p^e,\quad 0\le s<e,
 \label{sp:eq:affine-relation}
\end{equation}
where $U$ is homogeneous, $V$ is a coefficient-one monomial of the
same degree $d\ge1$, and
\[
 (d+1)p^s\le Q,\qquad a_i\in\Fp,\qquad
 1+\sum_i a_i=0,\qquad \beta\in\Fp^\times.
\]
None of these data is given. A deterministic polynomial sparse-bit
algorithm finds and verifies a representation whenever one exists,
normalizing $U$ so that $V$ is absent from its support, and replaces
$F$ by the nonzero rows among $U,V$. Zero normalized $U$ is allowed.
The generated ideals satisfy
\[
 K\subseteq K'\subseteq I,\qquad (K')^{\mathrm{sat}}=I.
\]
The total output support is at most the support of $F$, no coordinate
index is introduced, and every output row has degree $\deg(F)/Q$.
The algorithm can minimize this degree among its passing candidates.
Its repeated use with the preceding support-nonincreasing operations
has polynomial bit cost, including when the number of rows grows.
\end{theorem}
\begin{proof}
The change $U\mapsto U-cV$, $c\in\Fp$, changes the right side of
\eqref{sp:eq:affine-relation} by
$-c(1+\sum_i a_i)V^Q=0$. Subtracting the coefficient of $V$ therefore
gives the stated normalization without changing any other parameter.

As in Theorem~\ref{sp:thm:linearized-preprocessing}, a lower-block
exponent $Qv+p^i(u-v)$ is coordinatewise divisible by $Q$ only when
$u=v$: the coordinates of $u-v$ have absolute value at most $d$, while
$Q/p^i\ge d+1$. Thus the divided pure $Q$-section is exactly
\[
 W=E_Q(F)=U+\beta V.
\]
Its support is the disjoint union of that of $U$ and the monomial $V$.
For each power $Q$ dividing $D=\deg F$, and each $s$ satisfying the
gap, enumerate $V$ from the support of $W$. Its coefficient is the
only possible $\beta$, and deleting that term gives $U$. Solve the
sparse coefficient system
\[
 \sum_{i=0}^s a_i V^{Q-p^i}U^{p^i}
       =F-U^Q-\beta V^Q,
 \qquad \sum_{i=0}^s a_i=-1.
\]
There are $O(\log D)$ unknowns and $O(|\operatorname{supp}F|\log D)$
coefficient equations. Deterministic elimination and an exact sparse
identity check accept precisely the passing representations. The
true normalized representation is among them. If $U=0$, all block
columns are zero; the system accepts exactly when $F=\beta V^Q$,
and the sum constraint still has a solution. Ordering $Q$
decreasingly and fixing the monomial and pivot orders gives the
minimum output degree with deterministic ties.

For safety, suppose that the image of $V$ in $S/I$ were nonzero.
Then $t=U/V$ in its fraction field would be algebraic over $\Fp$.
Geometric integrality and the constant-field argument in
Theorem~\ref{sp:thm:linearized-preprocessing} give $t\in\Fp$.
The relation becomes
\[
 0=t^Q+\sum_i a_it^{p^i}+\beta
   =\left(1+\sum_i a_i\right)t+\beta=\beta,
\]
a contradiction. Hence $V\in I$, and the original relation gives
$U^Q\in I$, so $U\in I$. Every term of $F$ belongs to $(U,V)$.
The ideal sandwich and equality of saturations follow.

The section identity gives
\[
 |\operatorname{supp}U|+|\operatorname{supp}V|
   =|\operatorname{supp}W|\le|\operatorname{supp}F|.
\]
All output exponents are original section exponents divided by $Q$,
so no coordinate index is added. The degree decreases by a factor
$Q\ge p$. Every candidate block is obtained by a Frobenius power and
a monomial shift; neither large numerical powers nor coefficient
vectors are enumerated. This proves the cost of one recognition pass.

For iteration, let $T_0$ be the original total support, $D_0$ the
maximum original degree, and $v$ the number of occurring coordinates.
Every intermediate nonzero list has at most $T_0$ rows. Use
\[
 \Phi=\sum_{H\text{ a nonzero row}}
      |\operatorname{supp}H|\bigl(1+\log\deg H\bigr).
\]
An affine step decreases $\Phi$ by at least $\log p$, because its
total output support does not increase and both outputs have the
same smaller degree. A nontrivial rank-one step decreases it by at
least $1/(v+2)$, using the degree estimate of
Lemma~\ref{sci:lem:deflation}; the other preceding degree reductions
have the same or a stronger bound. Thus the number of these steps
is polynomial in $T_0,v,\log(D_0+1)$. Full eligible coordinate content
is removed in maximal batches between them. Each batch has at most
$v$ removals per row, since changing one coordinate does not alter
the other coordinates' minimum exponents. This also has polynomial
cost. Current supports, coefficient sizes, coordinate labels and
exponent lengths stay bounded by the original input parameters, so
all intervening recognition scans are polynomial.
\end{proof}

The conclusion forces a coordinate variable into $I$, since the
positive-degree monomial $V$ belongs to the prime ideal. It therefore
does not give an additional reduction on a projective scheme lying
in no coordinate hyperplane. A homogeneous relation with
$1+\sum_i a_i=0$ and $\beta=0$ needs a different output: the relation
allows every prime-field value of $U/V$, and does not force $U$ or
$V$ separately into $I$.
